% Part: sets-functions-relations
% Chapter: arithmetization
% Section: rationals

\documentclass[../../../include/open-logic-section]{subfiles}

\begin{document}

\olfileid{sfr}{arith}{rat}
\olsection{Van $\Int$ naar $\Rat$}
We hebben zojuist gezien hoe de gehele getallen met enige na\"{i}eve
verzamelingenleer uit de natuurlijke getallen kunnen worden geconstrueerd.
Nu construeren we de rationale getallen op vrijwel dezelfde wijze uit de
gehele getallen. Onze eerste observaties zijn:
\begin{enumerate}
	\item Elk rationaal getal kan worden geschreven in de vorm $\nicefrac{i}{j}$,
	waarbij $i$ en $j$ gehele getallen zijn en $j$ niet nul is.
	\item De informatie in een uitdrukking $\nicefrac{i}{j}$ kan evengoed
	worden vastgelegd met een geordend paar $\tuple{ i, j}$.
\end{enumerate}
Het ligt daarom voor de hand de rationale getallen op te vatten \emph{als}
geordende paren uit $\Int \times (\Int \setminus \{0_\Int\})$.
Net als eerder zou dat echter iets te na\"ief zijn: we willen immers
$\nicefrac{3}{2} = \nicefrac{6}{4}$, terwijl $\tuple{ 3, 2}\neq \tuple{ 6,
4}$. Algemener willen we:
\[
	\nicefrac{a}{b} = \nicefrac{c}{d} \text{ precies dan als } a \times d = b \times c
\]
Daartoe definiëren we als volgt een {equivalentierelatie} op $\Int \times
(\Int \setminus \{0_\Int\})$:
\[
	\tuple{ a, b }\Ratequiv \tuple{ c, d} \text{ precies dan als }a \times d = b \times c
\]
We moeten controleren dat dit inderdaad een equivalentierelatie is. Dat gaat
nagenoeg hetzelfde als bij $\Intequiv$ en laten we als oefening over.
\begin{prob}
Bewijs dat $\Ratequiv$ een equivalentierelatie is.
\end{prob}
Daarmee kunnen we de rationale getallen definiëren:
\begin{defn}
De rationale getallen zijn de equivalentieklassen onder $\Ratequiv$ van paren
van gehele getallen waarvan het tweede element niet nul is. Dat wil zeggen:
$\Rat =
\equivclass{(\Int \times (\Int\setminus \{0_\Int\}))}{\Ratequiv}$.
\end{defn}
Net als voor de gehele getallen willen we ook enkele basisbewerkingen
definiëren. Als $\equivrep{i,j}{\Ratequiv}$ de equivalentieklasse onder
$\Ratequiv$ is die $\tuple{i, j}$ als !!a{element} bevat, definiëren we:
\begin{align*}
	\equivrep{a, b}{\Ratequiv} + \equivrep{c, d}{\Ratequiv} &= \equivrep{ad + bc,  bd}{\Ratequiv}\\
	\equivrep{a, b}{\Ratequiv} \times \equivrep{c, d}{\Ratequiv} &= \equivrep{a  c, b d}{\Ratequiv}.
\intertext{Voor deze rationale getallen definiëren we $r \leq s$ met behulp
van het feit dat $r \le s$ precies dan als $s - r$ niet negatief is, dat wil
zeggen dat $s - r$ kan worden geschreven als $\nicefrac{i}{j}$ met $i$ niet-negatief en
$j$~positief:}
	\equivrep{a, b}{\Ratequiv} \leq \equivrep{c, d}{\Ratequiv} &\text{ precies dan als }
	\equivrep{c, d}{\Ratequiv} - \equivrep{a, b}{\Ratequiv} =
	\equivrep{i_\Int, j_\Int}{\Ratequiv}
\end{align*}
voor zekere $i \in \Nat$ en $0 \neq j \in \Nat$.
Vervolgens moeten we controleren of deze definities doen wat ze
\emph{behoren} te doen; die controle stellen we uit tot \olref[check]{sec}.
Ze doen dat inderdaad!{} Ten slotte willen we gehele getallen \emph{als}
rationale getallen kunnen behandelen. Daarom spreken we voor iedere
$i \in \Int$ af dat $i_\Rat = \equivrep{i, 1_\Int}{\Ratequiv}$. Ook hiervoor
gaan we in \olref[check]{sec} na dat alles naar behoren werkt.
\begin{prob}
Bewijs dat $(i + j)_\Rat = i_\Rat+ j_\Rat$ en $(i \times j)_\Rat =
i_\Rat \times j_\Rat$ en $i \leq j \liff i_\Rat \leq j_\Rat$, voor alle
$i, j \in \Int$.
\end{prob}
\end{document}
